\section*{ТЕРМИНЫ И ОПРЕДЕЛЕНИЯ}
\addcontentsline{toc}{section}{ТЕРМИНЫ И ОПРЕДЕЛЕНИЯ}

\begin{description}[leftmargin=0pt, style=nextline]

\item[Несанкционированные действия в сети]
действия, нарушающие установленную политику безопасности
корпоративной информационной системы, в том числе попытки получения
неавторизованного доступа, эксфильтрации данных, исчерпания ресурсов,
скрытого латерального перемещения и развёртывания вредоносного кода.

\item[Аномальная сетевая активность]
поведение сетевых соединений, статистически отклоняющееся от
характерного для контролируемой инфраструктуры профиля нормы,
проявляющееся в признаках длительности, направления, объёма и
последовательности соединений.

\item[Сетевой поток (flow)]
агрегированная единица сетевого трафика, описываемая совокупностью
характеристик \emph{(src/dst IP, src/dst port, protocol)} в пределах
заданного временного интервала и характеризуемая длительностью,
числом переданных пакетов и байтов, межпакетными интервалами и
сопутствующими статистиками.

\item[NIDS (Network-based Intrusion Detection System)]
сетевая система обнаружения вторжений, осуществляющая анализ
проходящего трафика для выявления попыток несанкционированного
доступа.

\item[NTA (Network Traffic Analysis)]
подкласс систем мониторинга сетевой безопасности, ориентированный
на анализ метаданных сетевых соединений (\emph{flow-level})
с применением статистических и машинно-обучаемых моделей.

\item[Аномальный детектор]
программный модуль, формирующий непрерывный скоринг
\(s(x)\colon X \to \mathbb{R}\) и принимающий решение об
аномальности наблюдения \(x\) при превышении порога \(\tau\).

\item[Скользящее окно (sliding window)]
последовательность из \(W\) подряд идущих flow-записей,
используемая в качестве входа модели; параметры окна --- размер
\(W\) и шаг \(S\).

\item[Дрейф распределений (concept/covariate drift)]
изменение совместного распределения \(p(x, y)\) или его компонент
во времени в продуктивной эксплуатации; источник деградации
аномального детектора.

\item[Калибровка вероятностей]
процедура согласования выходной вероятности классификатора
с эмпирической частотой положительного класса; в работе
используется температурное масштабирование.

\item[ИИ-злоумышленник (AI-attacker)]
активный программный компонент, формирующий поток сетевых событий
с управляемым переключением режимов \emph{норма/атака/дрейф} в
рамках программного стенда; в данной работе реализован на основе
библиотеки параметризованных шаблонов, конечного автомата
планирования и инжектора дрейфа.

\item[SOC (Security Operations Center)]
функциональная организационно-техническая структура мониторинга
и реагирования на инциденты информационной безопасности
корпоративного класса.

\item[Алерт]
сформированное системой обнаружения событие безопасности,
сопровождаемое уровнем критичности, скорингом, типом
предположительной атаки и контекстом наблюдения, направляемое в
SOC для последующей обработки аналитиком.

\end{description}
